\chapter{Referencial teórico}

Este capítulo apresenta os pilares teóricos que alicerçam a arquitetura e o desenvolvimento do sistema de API proposto, sendo dividido em três tópicos complementares: primeiramente, os princípios da Orientação a Objetos (OO), que trazem a base conceitual sobre a qual todo o desenvolvimento é edificado; em seguida, são abordados os princípios SOLID, que fortalecem as diretrizes para o \textit{design} das classes e módulos, tornando-os mais flexíveis e manuteníveis; e, por fim, é apresentado o \textit{Domain-Driven Design} (DDD), que orienta o desenvolvimento e o entendimento do sistema de forma alinhada ao domínio do negócio. A decisão por apresentá-los nessa ordem não é aleatória: cada tema é tratado como uma camada que complementa a anterior, formando um conjunto de conhecimentos coeso e aplicável a domínios com regras de negócio complexas e relativamente estáveis, como é o caso da doação sanguínea.

\section{Orientação a Objetos}
\label{sec:oo}

A OO é um paradigma de programação que tem como objetivo organizar o código em torno de objetos, estruturas que encapsulam dados e comportamentos, simulando entidades e eventos do mundo real \cite{sommerville2011}. Embora tenha surgido nos anos 1960 com a linguagem de programação Simula, a OO só ganhou força mais tarde, com linguagens como Smalltalk, C++ e Java, estabelecendo-se como paradigma fundamental no desenvolvimento de \textit{software} \cite{booch2007}. A disseminação desse paradigma na comunidade de desenvolvedores decorre de sua capacidade de modelar problemas do mundo real de forma intuitiva, aproximando a estrutura do código e a organização das classes da estrutura do problema que se pretende resolver \cite{sommerville2011}.

A OO é estruturada em quatro pilares fundamentais e complementares: abstração, encapsulamento, herança e polimorfismo.

\subsection{Abstração}

A abstração permite que os programadores se preocupem apenas com as características essenciais de um objeto, ignorando detalhes desnecessários ao modelo e ao contexto em questão \cite{booch2007}. Em vez de lidar com cada detalhe, desenvolvedores trabalham em níveis mais altos de abstração. Isso ocorre porque, segundo \citeonline{sommerville2011}, a abstração ajuda a controlar a complexidade do \textit{software}, mantendo apenas o essencial do domínio na visão dos programadores.

A definição de contratos entre tipos surge principalmente por meio de classes abstratas e interfaces, que estabelecem obrigações funcionais sem detalhar o processo interno \cite{gamma1994}. Implementações concretas surgem depois, moldadas conforme as particularidades de cada classe descendente.

Na prática deste projeto, essa abstração reflete-se na criação da entidade genérica de participante do sistema, reunindo características comuns a diferentes atores --- como nome, contato e endereço --- independentemente de se tratar de um doador, de um solicitante ou de uma instituição de saúde. Especializações dessa abstração assumem comportamentos e regras próprias de cada papel, sem que o restante da aplicação precise conhecer particularidades internas para operar em alto nível.

\subsection{Encapsulamento}

O encapsulamento é responsável por controlar o acesso aos dados internos de um objeto, exigindo que qualquer interação aconteça por meio de vias específicas, o que preserva a coerência do estado interno \cite{sommerville2011}. Ainda segundo \citeonline{booch2007}, esconder os detalhes de um objeto que não contribuem para suas características essenciais faz parte do conceito, de maneira que a interface pública e a implementação interna permaneçam separadas, permitindo alterações na implementação que não afetam o código cliente que utiliza a classe \cite{sommerville2011}.

De maneira prática, o encapsulamento usa modificadores de acesso para definir quais atributos e métodos podem ser vistos ou utilizados externamente. Assim, os dados ficam protegidos, enquanto métodos específicos permitem interação controlada. De acordo com \citeonline{bloch2018}, essa prática permite validações e garante a consistência do objeto.

Para o cenário hemoterápico, a aplicação desse princípio assegura a proteção direta de regras críticas relativas à elegibilidade do doador e à validade das solicitações de sangue: informações como tipo sanguíneo, data da solicitação e prazo limite não sofrem alterações arbitrárias externas, sendo modificadas exclusivamente por operações que validam a consistência do estado (como garantir que o prazo limite seja posterior à data de abertura).

\subsection{Herança}

O conceito de herança possibilita o reaproveitamento e a especialização de comportamentos por subclasses, sendo fundamental para a reutilização de código e o estabelecimento de relações hierárquicas entre tipos \cite{sommerville2011}. Segundo \citeonline{booch2007}, a herança funciona como o meio pelo qual uma classe herda métodos e propriedades de outra, concretizando vínculos do tipo ``é-um'' entre entidades. Dessa forma, surgem arranjos hierárquicos que organizam classes representando relacionamentos ontológicos do contexto do problema.

A herança pode ocorrer de duas principais maneiras: a herança simples, em que uma classe deriva de uma única superclasse, e a herança múltipla, em que uma classe herda de mais de uma superclasse, recebendo todos os atributos e métodos não privados das classes-base. Em ambos os casos é possível sobrescrever métodos para implementação específica \cite{booch2007}. \citeonline{sommerville2011} destaca que a herança facilita a manutenção do sistema, uma vez que mudanças na superclasse são automaticamente propagadas para as subclasses.

Ao modelar os diferentes atores envolvidos, a herança permite derivar especializações (como pessoas físicas e instituições) a partir de uma estrutura comum de participante, concentrando atributos compartilhados na classe base e evitando a duplicação de regras cadastrais.

\subsection{Polimorfismo}

Formas distintas de um mesmo comportamento surgem quando objetos diversos utilizam uma interface idêntica, conforme descrito por \citeonline{booch2007}. Segundo \citeonline{sommerville2011}, o polimorfismo permite que métodos funcionem de formas diversas em diferentes níveis de abstração, característica essencial na programação orientada a objetos.

Existem dois tipos principais de polimorfismo: o polimorfismo de sobrecarga (\textit{overloading}), em que operações distintas com o mesmo nome coexistem com assinaturas diferentes, e o polimorfismo de sobrescrita (\textit{overriding}), em que subclasses redefinem métodos herdados da superclasse \cite{booch2007}. O polimorfismo de subtipo, também conhecido como polimorfismo de inclusão, permite que uma referência de uma superclasse ou interface aponte para objetos de qualquer uma de suas subclasses, possibilitando que o sistema determine em tempo de execução qual implementação específica deve ser invocada \cite{sommerville2011}. Essa característica reduz o acoplamento entre componentes e facilita a adição de novos tipos ao sistema sem modificar código existente. Para que o polimorfismo de subtipo funcione corretamente, é fundamental que as classes derivadas possam substituir suas classes mãe sem alterar o comportamento esperado do sistema, princípio conhecido como \textit{Liskov Substitution Principle} (LSP), segundo o qual objetos de uma classe derivada devem poder ser usados no lugar de objetos da classe base sem comprometer a correção do programa \cite{martin2003}.

Em uma API voltada à captação, o polimorfismo viabiliza tratar múltiplos perfis de uso por meio de interfaces unificadas, permitindo que a aplicação execute rotinas genéricas --- tais como envio de notificações ou verificação de autorização --- sem necessitar de condicionais explícitas sobre o tipo concreto de participante envolvido.

\section{Princípios SOLID}
\label{sec:solid}

Com origem na década de 2000, SOLID --- acrônimo para \textit{Single Responsibility}, \textit{Open-Closed}, \textit{Liskov Substitution}, \textit{Interface Segregation} e \textit{Dependency Inversion} --- reúne cinco princípios fundamentais voltados ao \textit{design} orientado a objetos em desenvolvimento de \textit{software}, originalmente pensados por Robert C. Martin \citeyear{martin2003}. Posteriormente, em 2004, Michael Feathers organizou essas contribuições na abreviação que tornou o conjunto mais acessível pelo acrônimo SOLID.

Partindo dos conceitos definidos por Martin, surgem diretrizes capazes de tornar os sistemas mais estruturados, desacoplados conforme a necessidade e mais simples de manter \cite{martin2008}. Embora nem sempre aplicado integralmente, segundo \citeonline{martin2017}, o uso rigoroso dessas diretrizes faz com que o \textit{software} cresça com maior controle, reduzindo consequências indesejadas ao modificar funcionalidades ou responder a novas demandas.

A fundação das diretrizes SOLID está ligada à maneira como elas lidam com falhas típicas no processo de criação de \textit{softwares}, como código rígido, frágil ou pouco reaproveitável \cite{martin2003}. Embora cada princípio trate de um aspecto distinto do projeto orientado a objetos, todos convergem para objetivos comuns: diminuir dependências entre partes distintas, fortalecer a coesão interna das classes e permitir expansões futuras sem alterar o que já está pronto. Em domínios com lógica de negócio complexa, como o da doação sanguínea, seguir tais princípios torna-se especialmente relevante.

Abaixo, os cinco princípios supracitados são detalhados um a um.

\subsection{\textit{Single Responsibility Principle} (SRP)}

Este princípio estabelece que uma classe deve possuir apenas uma única razão para mudar, concentrando-se em uma responsabilidade coesa \cite{martin2003}. Como destaca \citeonline{booch2007}, a separação de responsabilidades é essencial para gerir a ``complexidade organizada'', garantindo que alterações em uma regra de negócio não propaguem erros inesperados em partes não relacionadas do sistema.

\subsection{\textit{Open/Closed Principle} (OCP)}

Preconiza que as entidades de \textit{software} (classes, módulos, funções) devem estar abertas para extensão, mas fechadas para modificação \cite{martin2003}. O objetivo é permitir que o comportamento de um sistema seja estendido sem a necessidade de alterar o código-fonte original, o que é alcançado por meio do uso de abstrações e polimorfismo. Padrões de projeto como \textit{Strategy} e \textit{Decorator} são exemplificações práticas deste princípio \cite{gamma1994}.

\subsection{\textit{Liskov Substitution Principle} (LSP)}

Formulado por Barbara Liskov, este princípio define que objetos de uma classe derivada devem ser capazes de substituir objetos da classe base sem comprometer a semântica do programa \cite{martin2003}. A conformidade com o LSP garante que a hierarquia de herança seja logicamente consistente, assegurando que o polimorfismo de subtipo não resulte em comportamentos inesperados.

\subsection{\textit{Interface Segregation Principle} (ISP)}

Orienta que os clientes não devem ser forçados a depender de interfaces que não utilizam \cite{martin2003}. Em termos práticos, é preferível criar interfaces específicas e granulares em vez de uma interface única e sobrecarregada. Essa prática minimiza o impacto de mudanças e respeita o conceito de ocultação de informação e encapsulamento \cite{booch2007}.

\subsection{\textit{Dependency Inversion Principle} (DIP)}

Define que módulos de alto nível não devem depender de módulos de baixo nível, mas ambos devem depender de abstrações \cite{martin2003}. O fundamento deste princípio é ``programar para uma interface, não para uma implementação'', reduzindo o acoplamento e facilitando a injeção de dependências, o que aumenta a flexibilidade e a testabilidade do sistema \cite{martin2008}.

\section{\textit{Domain-Driven Design} (DDD)}
\label{sec:ddd}

Proposto originalmente por \citeonline{evans2003}, o DDD --- também conhecido como Projeto Orientado ao Domínio --- é uma abordagem de desenvolvimento de \textit{software} voltada para sistemas complexos cuja principal dificuldade reside na lógica de negócio. Por isso, tem como foco central o entendimento profundo das regras de negócio, estabelecendo que o \textit{design} do \textit{software} deve ser guiado por um modelo que reflita fielmente o domínio --- o problema central que a aplicação visa resolver.

Segundo \citeonline{vernon2013}, o DDD não deve ser encarado como uma tecnologia ou metodologia rígida, mas sim como um conjunto de ferramentas que permite lidar com a complexidade no coração do \textit{software}, priorizando o que é estrategicamente mais importante para a organização.

A base fundamental do DDD está na colaboração estreita entre desenvolvedores e especialistas do domínio --- aqueles que conhecem todas as nuances do negócio. Para que esse diálogo seja possível e produtivo, o DDD propõe a criação de uma Linguagem Ubíqua: um vocabulário comum e rigoroso, compartilhado por todos os membros da equipe, possibilitando um contato produtivo entre desenvolvedores, que a princípio não conhecem o domínio, e especialistas, que são a fonte de conhecimento para o desenvolvimento. Conforme \citeonline{vernon2016}, essa linguagem não deve existir apenas em documentos ou diagramas, mas precisa ser falada nas reuniões do dia a dia e refletida diretamente no código-fonte, de modo que a intenção do negócio nunca se perca nas decisões técnicas. No domínio da doação sanguínea, essa linguagem se traduziria na adoção de termos como ``doador'', ``solicitação de doação'' e ``tipo sanguíneo'' diretamente no vocabulário técnico da equipe, nomes que qualquer especialista da área reconheceria imediatamente, sem necessidade de tradução.

Para organizar esse conhecimento e tornar a abordagem mais acessível na prática, autores como Vernon e Khan estruturam o DDD em três pilares fundamentais --- e é justamente essa divisão que orienta as seções seguintes.

\subsection{Padrões de \textit{design} estratégico}

No DDD, o \textit{design} estratégico fornece diretrizes para a organização do sistema em larga escala, permitindo a identificação de fronteiras claras e o alinhamento do \textit{software} com os objetivos de negócio. O conceito central nesta fase é o Contexto Delimitado (\textit{Bounded Context}), que atua como uma fronteira semântica dentro da qual um modelo específico é aplicável e consistente \cite{evans2003}. A definição dessas fronteiras é essencial para evitar que o sistema degenere em uma estrutura técnica confusa e fortemente acoplada, frequentemente descrita na literatura como uma ``Grande Bola de Lama'' (\textit{Big Ball of Mud}) \cite{vernon2013}.

Para gerenciar a comunicação entre diferentes contextos, utiliza-se o Mapeamento de Contextos (\textit{Context Mapping}), que define as relações técnicas e organizacionais entre as partes do sistema \cite{vernon2016}. Dentre os padrões de mapeamento, destaca-se a Camada Anticorrupção (\textit{Anticorruption Layer}), que funciona como um mecanismo defensivo para isolar o modelo de domínio principal de influências externas ou sistemas legados, preservando a integridade das regras internas \cite{evans2003}.

No domínio da doação sanguínea, é comum delimitar ao menos dois contextos com preocupações distintas: um voltado à solicitação e à priorização de doações, centrado nas regras de compatibilidade sanguínea e de urgência, e outro voltado à gestão dos doadores, responsável pelas regras de elegibilidade, histórico de doações e comunicação com os candidatos. Uma separação dessa natureza favorece que alterações nas regras de um dos contextos --- por exemplo, os critérios de elegibilidade de um doador --- não impactem diretamente a lógica do outro.

\subsection{Padrões de \textit{design} tático}

Enquanto o \textit{design} estratégico lida com as fronteiras de larga escala, os padrões de \textit{design} tático fornecem as ferramentas para a modelagem detalhada dentro de um Contexto Delimitado, garantindo que os objetos do sistema possuam comportamento rico e mantenham a consistência dos dados \cite{evans2003}.

O primeiro bloco de construção são as Entidades: objetos definidos por uma identidade única que persiste ao longo do tempo, independentemente de mudanças em seus atributos \cite{vernon2013}. No domínio da doação sanguínea, um solicitante ou um centro de coleta são exemplos de entidades --- um solicitante continua sendo o mesmo solicitante mesmo que seu contato ou endereço seja atualizado, pois o que o identifica é sua identidade no domínio, não seus dados.

Em contrapartida, os Objetos de Valor (\textit{Value Objects}) são imutáveis e definidos exclusivamente por seus atributos, sem identidade própria \cite{evans2003}. Conceitos como tipo sanguíneo, endereço ou contato se enquadrariam nessa categoria: dois doadores com o mesmo tipo sanguíneo compartilham o mesmo valor, e não faz sentido distingui-los por identidade --- o que importa é o valor em si. Essa imutabilidade traz segurança ao modelo, pois garante que tais conceitos não sejam alterados de forma inadvertida.

Os Agregados representam agrupamentos de entidades e objetos de valor tratados como uma unidade para fins de consistência e mudança de dados. Cada agregado possui uma Raiz do Agregado (\textit{Aggregate Root}), responsável por garantir que todas as regras de negócio e invariantes sejam satisfeitas antes de qualquer transação ser concluída \cite{vernon2013}. No domínio da doação sanguínea, uma solicitação de doação --- com seu ciclo de vida completo, desde a criação até o atendimento --- ou o histórico de doações de um doador são exemplos plausíveis de agregados: nenhum elemento interno poderia ser manipulado diretamente por objetos externos, devendo toda interação passar obrigatoriamente pela raiz do agregado, o que preserva a integridade do modelo.

Por fim, os Eventos de Domínio registram ocorrências significativas para o negócio no passado, permitindo a comunicação desacoplada entre diferentes partes do sistema \cite{vernon2016}. No contexto da doação sanguínea, a confirmação de uma doação ou a manifestação de interesse de um candidato são exemplos de eventos capazes de disparar notificações ou atualizar o estado de uma solicitação, sem que os componentes envolvidos precisem se conhecer diretamente.

\subsection{Arquitetura e isolamento do domínio}

Para que o modelo de domínio permaneça livre de contaminações tecnológicas, o DDD recomenda arquiteturas que isolem a lógica de negócio das preocupações de infraestrutura. Padrões como a Arquitetura em Camadas ou a Arquitetura de Portas e Adaptadores --- também conhecida como Arquitetura Hexagonal --- garantem que a camada de domínio permaneça pura e independente de detalhes técnicos, como bancos de dados ou interfaces de usuário \cite{evans2003}. Esse isolamento facilita a manutenibilidade, a testabilidade e a evolução contínua do conhecimento sobre o negócio dentro do \textit{software} \cite{vernon2016}.

Em um sistema modelado sob essa perspectiva, esse isolamento se refletiria na separação clara entre as classes de domínio --- responsáveis por expressar as regras de negócio da doação sanguínea --- e os serviços responsáveis por coordenar operações de infraestrutura, como persistência e envio de notificações. Dessa forma, a camada de domínio permaneceria como uma abstração independente de tecnologia, sobre a qual o restante da aplicação é construído.
